% 連続体力学 連続体力学-1 問題4（004-accel） 解いた日 2026-10-20
%% kotae
\toi{1}{○}{記述}
\[ \vect{a}=\frac{\mathrm{D}\vect{u}}{\mathrm{D}t}=\pdv{\vect{u}}{t}+(\vect{u}\cdot\vect{\nabla})\vect{u} \]
\nuke{$a_{j}=\mathrm{D}u_{j}/\mathrm{D}t$，式(4)で $A=u_{j}$}
\[ a_{j}=\pdv{u_{j}}{t}+\sum_{i=1}^{3}u_{i}\pdv{u_{j}}{x_{i}} \]
\naze{示すべき式をそのまま書いた形で，導出の段が無い。式(4)はスカラー量 $A$ についての関係なので，加速度の第 $j$ 成分 $a_{j}=\mathrm{D}u_{j}/\mathrm{D}t$ をとり，$u_{j}$ を1つのスカラー量とみなして $A=u_{j}$ とおく，と1行書けば導出になる。}
\toi{2}{◎}{}
\begin{align*}
\vect{a}&=\frac{\mathrm{D}\vect{u}}{\mathrm{D}t}\\
&=\pdv{\vect{u}}{t}-\Omega x_{2}\pdv{\vect{u}}{x_{1}}+\Omega x_{1}\pdv{\vect{u}}{x_{2}}\\
&=0-\Omega x_{2}(0,\Omega,0)+\Omega x_{1}(-\Omega,0,0)\\
&=(-\Omega^{2}x_{1},-\Omega^{2}x_{2},0)
\end{align*}
（図：点から原点の向きへの矢印）
\[ \abs{\vect{a}}=\Omega^{2}r \]
\toi{3}{◎}{}
\[ \vect{u}=\frac{\mathrm{D}\vect{x}}{\mathrm{D}t} \]
\[ -\Omega x_{2}=\frac{\mathrm{D}x_{1}}{\mathrm{D}t} \]
\[ \Omega x_{1}=\frac{\mathrm{D}x_{2}}{\mathrm{D}t},\quad \therefore\ \Omega\frac{\mathrm{D}x_{1}}{\mathrm{D}t}=\frac{\mathrm{D}^{2}x_{2}}{\mathrm{D}t^{2}} \]
\[ 0=\frac{\mathrm{D}x_{3}}{\mathrm{D}t},\quad \therefore\ x_{3}=X_{3} \]
\[ \frac{\mathrm{D}^{2}x_{2}}{\mathrm{D}t^{2}}=-\Omega^{2}x_{2} \]
\[ x_{2}(t)=C_{1}\sin(\Omega t)+C_{2}\cos(\Omega t) \]
\[ x_{2}(0)=C_{2}=X_{2} \]
\begin{align*}
x_{1}(t)&=\frac{1}{\Omega}\bigl(C_{1}\Omega\cos(\Omega t)-C_{2}\Omega\sin(\Omega t)\bigr)\\
&=C_{1}\cos(\Omega t)-C_{2}\sin(\Omega t)
\end{align*}
\[ x_{1}(0)=C_{1}=X_{1} \]
\[ \vect{x}(t)=\begin{pmatrix}X_{1}\cos(\Omega t)-X_{2}\sin(\Omega t)\\X_{1}\sin(\Omega t)+X_{2}\cos(\Omega t)\\X_{3}\end{pmatrix} \]
\begin{align*}
\frac{\mathrm{D}^{2}\vect{x}}{\mathrm{D}t^{2}}&=\begin{pmatrix}-\Omega^{2}\bigl(X_{1}\cos(\Omega t)-X_{2}\sin(\Omega t)\bigr)\\-\Omega^{2}\bigl(X_{1}\sin(\Omega t)+X_{2}\cos(\Omega t)\bigr)\\0\end{pmatrix}\\
&=\begin{pmatrix}-\Omega^{2}x_{1}\\-\Omega^{2}x_{2}\\0\end{pmatrix}
\end{align*}
\toi{4}{○}{記述}
\begin{align*}
\vect{a}&=\frac{\mathrm{D}\vect{u}}{\mathrm{D}t}=\pdv{\vect{u}}{t}+\frac{x_{1}}{t+t_{0}}\pdv{\vect{u}}{x_{1}}\\
&=-\frac{x_{1}}{(t+t_{0})^{2}}+\frac{x_{1}}{t+t_{0}}\cdot\frac{1}{t+t_{0}}\\
&=0
\end{align*}
\[ \vect{u}=\frac{\mathrm{D}\vect{x}}{\mathrm{D}t},\quad \frac{x_{1}}{t+t_{0}}=\frac{\mathrm{D}x_{1}}{\mathrm{D}t} \]
\[ \therefore\ \frac{\mathrm{D}t}{t+t_{0}}=\frac{\mathrm{D}x_{1}}{x_{1}} \]
\[ \log(t+t_{0})+C=\log x_{1} \]
\[ \therefore\ x_{1}(t)=C(t+t_{0}) \]
$\vect{x}(0)=(X_{1},X_{2},X_{3})$ とすると
\[ \vect{x}(t)=\begin{pmatrix}\frac{X_{1}}{t_{0}}(t+t_{0})\\X_{2}\\X_{3}\end{pmatrix} \]
固定点では粒子は加速しているように見えるが、流れる粒子は等速直線運動をしているから。\nuke{粒子の入れ替わり}
\naze{各粒子が等速であることは $\vect{a}=\vect{0}$ の理由になっているが，固定点で速度が変わる（$\partial\vect{u}/\partial t\neq\vect{0}$）理由が無い。速度 $X_{1}/t_{0}$ の小さい（出発点が原点に近い）粒子ほど遅れてその点に着き，次々と入れ替わるので，固定点の速度 $x_{1}/(t+t_{0})$ は減っていく。「加速しているように見える」より「速度が変わる（減る）」が正確である。}
%% kizuki
\item \mae{第1回 問題4・10-17}(1) で前回と同じく「$A=u_{j}$ とおく」段が抜けた。式(4)はスカラーの式なので，ベクトルに使うときは「第 $j$ 成分 $a_{j}=\mathrm{D}u_{j}/\mathrm{D}t$ で $A=u_{j}$」と最初に1行書いてから成分の式を書く。
\item \mae{第1回 問題4・10-17}(4) の理由も前回と同じ所が抜けた。粒子ごとに見た量（速度 $X_{1}/t_{0}$）が一定でも，固定点には別の粒子が次々と来るので場の値は変わる。軌道 $x_{1}=X_{1}(t+t_{0})/t_{0}$ から，時刻 $t$ に $x_{1}$ に来る粒子の速度 $X_{1}/t_{0}=x_{1}/(t+t_{0})$ が後ほど小さいことを1文で書く（問題1 (4) と同じ形）。
